\uuid{GJea}
\titre{Résolution d'équation différentielle du premier ordre}
\auteur{bodin ; croizat ; sacre}
\organisation{exo7}
\licence{CC BY-NC-SA 4.0}
\datecreate{2022-03-08}
\niveau{L1}
\module{Analyse}
\chapitre{Equation différentielle}
\sousChapitre{Résolution d'équation différentielle du premier ordre}
\difficulte{3}
\theme{}
\exo7qcmid{hilisit.equations-differentielles.052}

\contenu{
\begin{question}
\qclassification{225.01}
Soit $f$ une fonction dont la courbe représentative admet pour tangente en $x=-1$ la droite d'équation $y=2x-2$. Parmi les équations différentielles suivantes, quelle est la seule dont $f$ peut être une solution ?
\begin{answers}
  \bad{$y'=y+ e^x$}
  \good{$y'=-y+2x$}
  \bad{$y'=2y+3x^3$}
  \bad{$2y'-y=2$}
\end{answers}
\begin{explanations}
Sur la droite $y=2x-2$, le point d'abscisse $x=-1$ est $A(-1,-4)$ donc $f(-1)=-4$. De plus, la pente de la droite est $2$, donc $f'(-1)=2$. Parmi les équations différentielles proposées, $y'=-y+2x$ est la seule qui permet d'obtenir ces deux valeurs (on remplace $y$ par $f$, $y'$ par $f'$, et on évalue tout cela en $x=-1$).
\end{explanations}
\end{question}
}
